\subsection{Simulation \& Modeling --- MCQs}

\begin{mcqbox}
\textbf{Q1.} A queue holds 0 customers for 45 minutes and 6 customers for 15 minutes. The time-average queue length is\\
(a) $3$ \quad (b) $1.5$ \quad (c) $6$ \quad (d) $0.75$\\
\textbf{Ans: (b)} --- $(0\times45 + 6\times15)/60 = 1.5$. Averaging the two \emph{values} to get 3 ignores how long each lasted.

\textbf{Q2.} Customers arrive at a rate of $12$/hour and a single server handles $15$/hour. The utilization is\\
(a) $1.25$ \quad (b) $0.8$ \quad (c) $0.2$ \quad (d) $180$\\
\textbf{Ans: (b)} --- $\rho = \lambda/\mu = 12/15 = 0.8$, comfortably inside the healthy $70$--$85\%$ band.

\textbf{Q3.} If $\lambda > \mu$ in a single-server queue, then as time grows\\
(a) $\hat u$ exceeds 1 \quad (b) the queue length grows without bound \quad (c) the system reaches steady state \quad (d) the delay approaches zero\\
\textbf{Ans: (b)} --- utilization saturates at 1 (it is a fraction of time); it is $Q(t)$ that explodes.

\textbf{Q4.} In a discrete-event simulation, state variables change\\
(a) continuously \quad (b) at fixed time steps \quad (c) only at events \quad (d) at every clock tick\\
\textbf{Ans: (c)}

\textbf{Q5.} Which of these is a \emph{continuous-time} (time-weighted) statistic?\\
(a) $\hat d(n)$, the average delay \quad (b) $\hat q(n)$, the time-average queue length\\
(c) the number of customers served \quad (d) the total delay\\
\textbf{Ans: (b)} --- $\hat q$ and $\hat u$ are areas over time; $\hat d$ is a plain mean over customers.

\textbf{Q6.} The area accumulator at an event must use\\
(a) the new state value \quad (b) the old state value \quad (c) the average of the two \quad (d) either, it does not matter\\
\textbf{Ans: (b)} --- the old value is what actually held during the interval just elapsed. Getting this wrong is named in the notes as the most common simulation bug.

\textbf{Q7.} When the queue empties and the server goes idle, the next departure time is set to $\infty$ in order to\\
(a) end the simulation \quad (b) prevent a ``phantom departure'' being scheduled for a non-existent customer\\
(c) save memory \quad (d) force the clock to advance by one unit\\
\textbf{Ans: (b)}

\textbf{Q8.} A simulation model passes all verification checks but its predictions disagree with the real system. What failed?\\
(a) verification \quad (b) validation \quad (c) calibration only \quad (d) nothing; models always disagree\\
\textbf{Ans: (b)} --- the code correctly implements a model whose assumptions are wrong.

\textbf{Q9.} In the 10-step simulation study, a failed validity check at Step 6 sends you back to\\
(a) Step 5 \quad (b) Step 4 \quad (c) Step 2 \quad (d) Step 1\\
\textbf{Ans: (c)} --- invalid output usually means the conceptual model or the data is wrong, not just the code.

\textbf{Q10.} Which statement is \textbf{FALSE}?\\
(a) Next-event time advance skips idle periods \quad (b) Simulated time equals real computer time\\
(c) The event list holds future events and their times \quad (d) Stochastic models produce random outputs\\
\textbf{Ans: (b)} --- 24 simulated hours may take 0.01 s, and 1 ns of a network may take 10 minutes.
\end{mcqbox}

\subsection{RNG \& Monte Carlo --- MCQs}

\begin{mcqbox}
\textbf{Q11.} For $R \sim U(0,1)$, $\operatorname{Var}(R)$ equals\\
(a) $1/2$ \quad (b) $1/3$ \quad (c) $1/12$ \quad (d) $1/4$\\
\textbf{Ans: (c)} --- $E[R^2]-E[R]^2 = \frac13-\frac14 = \frac1{12} \approx 0.0833$.

\textbf{Q12.} An LCG has $X_{i+1} = (5X_i+3)\bmod 16$ with $X_0 = 1$. Then $X_2$ is\\
(a) 8 \quad (b) 11 \quad (c) 10 \quad (d) 5\\
\textbf{Ans: (b)} --- $X_1 = 8$, $X_2 = (40+3)\bmod16 = 43\bmod16 = 11$.

\textbf{Q13.} Does that generator achieve the full period $m = 16$?\\
(a) no, because $m$ is not prime \quad (b) yes: $\gcd(3,16) = 1$ and $5 = 1+4(1)$\\
(c) no, because $c \neq 0$ \quad (d) only for even seeds\\
\textbf{Ans: (b)} --- both mixed-generator conditions hold, and the period really is 16.

\textbf{Q14.} For the multiplicative generator $X_{i+1} = 13X_i \bmod 64$, the maximum achievable period is\\
(a) 64 \quad (b) 63 \quad (c) 16 \quad (d) 32\\
\textbf{Ans: (c)} --- with $c=0$ and $m=2^b$ the ceiling is $m/4 = 16$, attained only for \emph{odd} seeds.

\textbf{Q15.} The values an LCG can produce after scaling are spaced\\
(a) $1/m$ apart \quad (b) $1/a$ apart \quad (c) continuously \quad (d) $1/c$ apart\\
\textbf{Ans: (a)} --- $R_i \in \{0, 1/m, \dots, (m-1)/m\}$, which is why a large $m$ matters.

\textbf{Q16.} A sequence gives $0.01, 0.02, 0.03, \dots, 0.99$. Which test will flag it?\\
(a) chi-square \quad (b) Kolmogorov--Smirnov \quad (c) autocorrelation \quad (d) none; it is a good generator\\
\textbf{Ans: (c)} --- it is perfectly uniform (so both uniformity tests pass) but utterly predictable. Only an independence test catches it.

\textbf{Q17.} In the autocorrelation test with $i=2$, $\ell=4$, $N=30$, the value of $M$ is\\
(a) 5 \quad (b) 6 \quad (c) 7 \quad (d) 4\\
\textbf{Ans: (b)} --- largest $M$ with $i+(M+1)\ell \le N$: $2+7(4) = 30 \le 30$ $\checkmark$, while $2+8(4) = 34 > 30$.

\textbf{Q18.} The $0.25$ subtracted in $\hat\rho_{i\ell}$ arises because\\
(a) it is an arbitrary correction \quad (b) $E[R]E[R'] = \frac12\cdot\frac12$ for independent uniforms\\
(c) $\operatorname{Var}(R) = 1/12$ \quad (d) there are 4 quadrants\\
\textbf{Ans: (b)} --- it centres the estimator at zero under independence.

\textbf{Q19.} An RNG is subjected to 20 independent tests at $\alpha = 0.05$ and fails 1. You should\\
(a) reject it immediately \quad (b) accept it: about $64\%$ of good generators fail at least one of 20 tests\\
(c) conclude the tests are broken \quad (d) rerun with $\alpha = 0.5$\\
\textbf{Ans: (b)} --- $1-0.95^{20} = 0.6415$, and the expected number of chance failures is $20(0.05) = 1$.

\textbf{Q20.} A Monte Carlo estimate must be made 4 times more accurate. The sample size must grow by a factor of\\
(a) 4 \quad (b) 8 \quad (c) 16 \quad (d) 2\\
\textbf{Ans: (c)} --- error $\propto 1/\sqrt n$, so $n$ scales as $4^2$.

\textbf{Q21.} In estimating $\pi$ by hit-or-miss, $\hat\pi = 4h/n$ because\\
(a) the circle has radius 4 \quad (b) $P(\text{hit}) = (\pi/4)/1$ for a quarter circle in the unit square\\
(c) there are 4 quadrants to sum \quad (d) $4$ normalizes the variance\\
\textbf{Ans: (b)}

\textbf{Q22.} Which statement is \textbf{FALSE}?\\
(a) Passing all tests does not prove a generator is random \quad (b) K-S works for small samples; chi-square needs large $N$\\
(c) A longer period guarantees independence \quad (d) The seed determines the entire sequence\\
\textbf{Ans: (c)} --- period length says nothing about correlation structure within the cycle.
\end{mcqbox}

\subsection{Simulation, RNG \& Monte Carlo --- Problems}

\begin{probbox}
\PQ{P1.} A single-server queue receives 4 customers with interarrival times $0.3,\ 0.9,\ 0.6,\ 1.2$ and service times $1.2,\ 0.8,\ 1.5,\ 0.5$. Build the trace and compute $\hat d(4)$, $\hat q(4)$, $\hat u(4)$ over the run ending when customer 4 departs. Then test Little's Law, and comment on stability.

\ans Arrival times are the cumulative sums: $0.3,\ 1.2,\ 1.8,\ 3.0$.
\begin{center}\small
\begin{tabular}{@{}ccccc@{}}
\toprule
Cust. & Arrives & Starts service & Departs & Delay $D_i$\\
\midrule
1 & 0.3 & 0.3 & 1.5 & 0.0\\
2 & 1.2 & 1.5 & 2.3 & 0.3\\
3 & 1.8 & 2.3 & 3.8 & 0.5\\
4 & 3.0 & 3.8 & 4.3 & 0.8\\
\bottomrule
\end{tabular}
\end{center}
(Each customer starts at $\max(\text{own arrival},\ \text{previous departure})$.) Run length $T = 4.3$.
\[
\hat d(4) = \frac{0+0.3+0.5+0.8}{4} = \frac{1.6}{4} = \mathbf{0.4}
\]
\emph{Queue area:} customer $i$ waits over $[\text{arrival},\text{start})$, contributing that length to $\int Q\,dt$:
\[
\int_0^{4.3}Q(t)\,dt = 0 + 0.3+0.5+0.8 = 1.6
\qquad\Rightarrow\qquad
\hat q(4) = \frac{1.6}{4.3} = \mathbf{0.372093}
\]
\emph{Busy area:} the server runs continuously from $0.3$ to $4.3$:
\[
\int_0^{4.3}B(t)\,dt = 4.0 \qquad\Rightarrow\qquad \hat u(4) = \frac{4.0}{4.3} = \mathbf{0.930233}
\]
\emph{Little's Law:} $\lambda = 4/4.3 = 0.930233$, $W_q = \hat d(4) = 0.4$, so
\[
\lambda W_q = 0.930233\times0.4 = 0.372093 = \hat q(4) \ \checkmark
\]
It holds \emph{exactly} here because every customer's queueing time was completed inside the window --- nobody was left waiting at $T$.

\emph{Stability:} $E[A] = 3.0/4 = 0.75$ so $\lambda = 1.333$; $E[S] = 4.0/4 = 1.0$ so $\mu = 1.0$; $\rho = 1.333 > 1$. The system is \textbf{unstable} --- work arrives faster than it clears, which is why each successive delay ($0, 0.3, 0.5, 0.8$) is larger than the last.

\PQ{P2.} For the LCG $X_{i+1} = (5X_i+3)\bmod16$ with $X_0 = 1$: generate 6 values and their $R_i$, and determine whether the full period $m=16$ is achieved.

\ans
\[
X_1 = (5+3)\bmod16 = 8,\quad
X_2 = (40+3)\bmod16 = 11,\quad
X_3 = (55+3)\bmod16 = 10
\]
\[
X_4 = (50+3)\bmod16 = 5,\quad
X_5 = (25+3)\bmod16 = 12,\quad
X_6 = (60+3)\bmod16 = 15
\]
\begin{center}\small
\begin{tabular}{@{}lcccccc@{}}
\toprule
$X_i$ & 8 & 11 & 10 & 5 & 12 & 15\\
$R_i = X_i/16$ & 0.5000 & 0.6875 & 0.6250 & 0.3125 & 0.7500 & 0.9375\\
\bottomrule
\end{tabular}
\end{center}
\emph{Max-period test} (mixed generator, $c\neq0$, $m = 2^4$): need (i) $\gcd(c,m) = 1$ and (ii) $a = 1+4k$.
\[
\gcd(3,16) = 1 \ \checkmark, \qquad a = 5 = 1+4(1) \ \checkmark
\]
Both hold, so the period is the full $\mathbf{16}$: every integer $0..15$ appears exactly once before repeating.

\emph{Contrast:} the generator $X_{i+1} = (3X_i+7)\bmod16$ has $\gcd(7,16)=1$ but $a = 3$ is \emph{not} of the form $1+4k$, and its period is only \textbf{8}.

\PQ{P3.} 120 random numbers are binned into 6 equal intervals with observed counts $18, 22, 25, 17, 20, 18$. Perform a chi-square uniformity test at $\alpha = 0.05$ (critical value $\chi^2_{0.05,5} = 11.07$).

\ans $E_i = N/n = 120/6 = 20$ for every bin.
\begin{center}\small
\begin{tabular}{@{}lcccccc|c@{}}
\toprule
$O_i$ & 18 & 22 & 25 & 17 & 20 & 18 & 120\\
$(O_i-E_i)^2$ & 4 & 4 & 25 & 9 & 0 & 4 & ---\\
$(O_i-E_i)^2/E_i$ & 0.2 & 0.2 & 1.25 & 0.45 & 0.0 & 0.2 & \textbf{2.30}\\
\bottomrule
\end{tabular}
\end{center}
\[
\chi_0^2 = 2.30, \qquad \text{df} = n-1 = 5, \qquad \chi^2_{0.05,5} = 11.07
\]
Since $2.30 < 11.07$, \textbf{do not reject $H_0$} --- no evidence against uniformity. Note this says nothing about \emph{independence}.

\PQ{P4.} A Monte Carlo estimator has $\sigma = 0.25$. (i) How many samples give a 95\% confidence half-width of at most $0.01$? (ii) How many for $0.0025$? (iii) State the general scaling law.

\ans \emph{(i)} Half-width $= 1.96\sigma/\sqrt n \le 0.01$:
\[
\sqrt n \ge \frac{1.96(0.25)}{0.01} = 49 \qquad\Rightarrow\qquad n \ge \mathbf{2401}
\]
\emph{(ii)} The target is $4\times$ tighter, so $\sqrt n$ must grow $4\times$ and $n$ by $4^2 = 16$:
\[
n \ge 16(2401) = \mathbf{38{,}416}
\]
(Direct check: $(1.96\times0.25/0.0025)^2 = 196^2 = 38416$ $\checkmark$.)

\emph{(iii)} Error $\propto \sigma/\sqrt n$, so \textbf{to improve accuracy by a factor $f$ you need $f^2$ times as many samples}. Ten times the accuracy costs a hundred times the work --- the fundamental limitation of Monte Carlo.

\PQ{P5.} A correct random-number generator is subjected to (i) 10 independent tests, (ii) 100 independent tests, all at $\alpha = 0.05$. In each case give the probability of at least one false rejection and the expected number of rejections. Then state the practical rule.

\ans Using $P(\ge1\text{ false reject}) = 1-(1-\alpha)^k$ and expected count $= k\alpha$:
\begin{center}\small
\begin{tabular}{@{}lcc@{}}
\toprule
$k$ & $1-0.95^k$ & expected rejections\\
\midrule
1 & 0.0500 & 0.05\\
5 & 0.2262 & 0.25\\
10 & \textbf{0.4013} & 0.50\\
20 & 0.6415 & 1.00\\
50 & 0.9231 & 2.50\\
100 & \textbf{0.9941} & 5.00\\
\bottomrule
\end{tabular}
\end{center}
\emph{(i)} 10 tests: about a \textbf{40\%} chance of at least one failure, expected $0.5$ failures.
\emph{(ii)} 100 tests: a \textbf{99.4\%} chance of at least one failure, expected \textbf{5} failures.

\emph{Practical rule:} a handful of failures among many tests is exactly what a \emph{good} generator produces. Worry only when the failure count is well above $k\alpha$, when the \emph{same} test fails repeatedly, or when several \emph{different} tests fail together.
\end{probbox}

\begin{trapbox}
\begin{itemize}
\item \textbf{Time-average $\ne$ average of the values.} Weight every value by how long it lasted (0 for 45 min and 6 for 15 min gives $1.5$, not $3$).
\item \textbf{$\hat u$ can never exceed 1.} An overloaded system shows $\hat u \to 1$ and a diverging $Q(t)$, not $\hat u > 1$.
\item \textbf{$\rho = \lambda/\mu$ uses \emph{rates}.} If the problem gives mean times, invert first: $\lambda = 1/E[A]$, $\mu = 1/E[S]$.
\item \textbf{Areas use the old state value} and width $(\text{clock}-\text{time of last event})$, in that order.
\item \textbf{$\hat q$ counts only those \emph{waiting}}, excluding the customer in service; $\hat u$ counts the server, not the queue.
\item \textbf{Little's Law is asymptotic.} It matches exactly only when the customers counted and the area measured cover the same population --- leftover waiters at time $T$ break it.
\item \textbf{Verification $\ne$ validation}, and both validation failures return to Step 2.
\item \textbf{Uniformity does not imply independence.} $0.01, 0.02, 0.03,\dots$ passes K-S and chi-square and is completely deterministic.
\item \textbf{$M$ in the autocorrelation test is the largest integer with $i+(M+1)\ell \le N$} --- off-by-one here wrecks both $\hat\rho$ and $\sigma$.
\item \textbf{Chi-square df is (bins $-1$)}, not the sample size and not the number of bins.
\item \textbf{A single test failure does not condemn a generator} --- with 10 tests a good one fails at least once about 40\% of the time.
\item \textbf{Monte Carlo accuracy costs quadratically.} ``Twice as accurate'' means four times the samples, not twice.
\item \textbf{The multiplicative case ($c=0$, $m=2^b$) caps out at $m/4$}, and only for odd seeds --- it can never reach the full $m$.
\end{itemize}
\end{trapbox}
